\chapter{Standard Tableaux and Branching}

The spanning set of all polytabloids from Result~\ref{res:A04} is
usually much larger than the dimension of the Specht module. The
standard polytabloids give a basis whose labels track restriction to
\(S_{n-1}\). The same labels yield an integer recurrence for the
dimension, which we solve by a determinant and then by hook lengths.
Our conventions for diagrams and tabloids are those of
Results~\ref{res:A02} and~\ref{res:A04}.

\section{The standard polytabloid basis}

A tableau \(t\) of shape \(\lambda\vdash n\) is \emph{standard} if its
labels \(1,\ldots,n\) increase strictly to the right along each row and
down each column. Call it \emph{column-increasing} if it satisfies only
the column condition. The empty diagram has a single tableau, which is
standard.

\begin{proposition}[Independence of standard polytabloids]\label{res:B01}
The vectors \(e_t\) for standard \(t\) are linearly independent in
\(M^\lambda\).
\end{proposition}

\begin{proof}
For a tabloid \(u\), let \(N_{k,m}(u)\) count the labels at most \(m\)
in its first \(k\) rows. These counts for all \(k,m\) determine \(u\):
the difference \(N_{k,m}(u)-N_{k,m-1}(u)\) tells whether the label
\(m\) is among the first \(k\) rows. Thus
\[
  u\preceq v\quad\Longleftrightarrow\quad
  N_{k,m}(u)\le N_{k,m}(v)\quad\text{for every } k,m
\]
is a partial order on tabloids of the given shape.

For a column-increasing \(t\), the entries of each column retained
in the first \(k\) rows are the smallest entries of that column in
the required number. Reordering any column cannot increase the count
of entries at most \(m\) in these positions. Hence
\(\{ct\}\preceq\{t\}\) for each \(c\in C_t\). The coefficient of
\(\{t\}\) in \(e_t\) is \(1\) by Result~\ref{res:A04}, and all other
terms are strictly smaller in the order just defined. Distinct
standard tableaux have distinct tabloids, since their row order can
be recovered by sorting each row.

If a finite linear combination of standard polytabloids had sum zero
and some coefficient were nonzero, take a maximal tabloid among the
leading tabloids with nonzero coefficient. No other summand contributes
at it, by the order property and distinctness. Its coefficient in the
sum is therefore nonzero, a contradiction. This is the triangular
argument of \cite[Section 2.5, Proposition 2.5.9]{sagan}.
\end{proof}

\begin{proposition}[Garnir relations and straightening]\label{res:B02}
Suppose \(t\) is column-increasing and has a row descent at adjacent
columns, with entries \(a_i>b_i\) in row \(i\). Write the two columns
as \(a_1<\cdots<a_p\) and \(b_1<\cdots<b_q\), and set
\[
  A=\{a_i,\ldots,a_p\},\qquad B=\{b_1,\ldots,b_i\}.
\]
For a set of labels \(U\), write \(S_U\) for the group that permutes
\(U\) and fixes all other labels. If \(P\) represents the left cosets
of \(S_A\times S_B\) in \(S_{A\cup B}\) and contains the identity, then
\[
  \sum_{\pi\in P}\operatorname{sgn}(\pi)e_{\pi t}=0.
\]
Every polytabloid is a linear combination of standard polytabloids.
\end{proposition}

\begin{proof}
Write \(U=A\cup B\) and \(H=S_A\times S_B\). Every term \(\{ct\}\)
of \(e_t\), with \(c\in C_t\), has the labels of \(U\) within the
same two columns. Those columns occupy at most \(p\) rows, since
\(q\le p\), while \(\lvert U\rvert=p+1\). Some two labels of \(U\)
therefore lie in the same row of \(ct\). Their transposition \(\tau\)
fixes \(\{ct\}\), and the terms indexed by \(\sigma\) and
\(\sigma\tau\) in the signed sum over \(S_U\) cancel on this tabloid.
The signed sum over \(S_U\) consequently annihilates \(e_t\).

The subgroup \(H\) is contained in \(C_t\), so \(h e_t\)
equals \(\operatorname{sgn}(h)e_t\) by Result~\ref{res:A04}. Factor the
signed sum over \(S_U\) using its disjoint cosets \(\pi H\). The inner
signed sum over \(H\) acts by \(\lvert H\rvert\), and dividing the
result by this nonzero integer proves the relation.

To straighten, first sort the columns of any tableau. A column
permutation acts on its polytabloid by its sign, so this step preserves
the span. Order the column-increasing tableaux by lexicographic order
on the word obtained by reading each column from top to bottom and
the columns from left to right. A column-increasing tableau that is
not standard has a descent as in the statement. If \(\pi\) represents
a nonidentity coset, then \(\pi A\ne A\). At least one label of \(A\)
is replaced in the left column by a label of \(B\). All labels of
\(B\) are smaller than every label of \(A\), because \(b_i<a_i\).
All other entries of this column are unchanged, and the earlier
columns are unchanged. After sorting, the first difference in the
left column is a smaller label, since the smallest label in the
symmetric difference of the two column sets belongs to the new set.
The column word has strictly decreased.

Isolate \(e_t\) in the relation and sort every other term, keeping the
sign of each sorting permutation. All resulting polytabloids have
smaller words, so induction over the finite order expresses each one
in the span of the standard polytabloids. The relation and
straightening are the column form of the argument in
\cite[Section 2.6, Proposition 2.6.3 and Theorem 2.6.4]{sagan}.
\end{proof}

\begin{theorem}[Standard basis and rational model]\label{res:B03}
The standard polytabloids form a basis of \(S^\lambda\). In
particular, \(\dim_{\mathbb C} S^\lambda=f^\lambda\), where
\(f^\lambda\) is the number of standard tableaux of shape \(\lambda\).
The span of all polytabloids in the rational tabloid module is a
rational \(S_n\)-module whose complexification is \(S^\lambda\).
\end{theorem}

\begin{proof}
Every generator of \(S^\lambda\) straightens by
Proposition~\ref{res:B02}, and Proposition~\ref{res:B01} gives
independence. The tabloid basis and every polytabloid have integral
coordinates. The signed coset relation and its use in straightening
hold over \(\mathbb Q\), and the same triangular proof gives
independence over \(\mathbb Q\). The resulting rational and complex
bases are indexed by the same standard tableaux. The action on
tabloids is given by permutations, so the scalar-extension map on
these bases is an isomorphism of \(S_n\)-modules
\cite[Theorem 2.6.5]{sagan}, \cite[Corollary 4.38]{etingof}.
\end{proof}

\section{Removing and adding one box}

A removable box of a nonempty diagram is a cell whose deletion
leaves a diagram. These are the last cells of rows \(r\) with
\(\lambda_r>\lambda_{r+1}\), where trailing parts are zero. For each
\(n\ge 1\), view \(S_{n-1}\) as the subgroup of \(S_n\) fixing the
label \(n\).

\begin{proposition}[Restriction filtration]\label{res:B04}
For \(n\ge 1\), let \(c_1,\ldots,c_s\) be the removable boxes of
\(\lambda\vdash n\), listed by increasing row numbers
\(r_1,\ldots,r_s\). Put
\(\mu_i=\lambda\setminus c_i\). For \(1\le i\le s\) let \(F_i\)
be the span of standard \(e_t\) with \(n\) in one of
\(c_1,\ldots,c_i\) and put \(F_0=0\). Then \(F_i\) are
\(S_{n-1}\)-submodules, and deleting \(n\) from the \(i\)-th row
choice induces \(F_i/F_{i-1}\simeq S^{\mu_i}\). Thus the full
restriction has a filtration with one successive quotient for
each removable box.
\end{proposition}

\begin{proof}
For each \(j\), define \(D_j:M^\lambda\to M^{\mu_j}\) on tabloids.
If \(n\) is in row \(r_j\), delete it from that row. Otherwise send
the tabloid to zero. This depends only on row sets and commutes with
all permutations fixing \(n\). The target row lengths are those of
\(\mu_j\) because \(c_j\) was removable. If this deletion empties the
last row, omit that row.

In a standard tableau \(t\), \(n\) has no neighbor to its right or
below, since such a neighbor would have a greater label. Suppose \(n\)
is in \(c_a\). In every term of \(e_t\), a column permutation can only
move \(n\) upwards within its column, since \(c_a\) is the bottom cell
of the column. Thus \(D_j(e_t)=0\) for \(j>a\). For \(j=a\), only
the column permutations fixing \(n\) contribute. Restricting them
to the other labels identifies them with the column group of
\(t\setminus n\) and preserves their signs. Therefore
\[
  D_a(e_t)=e_{t\setminus n}.
\]
Deleting \(n\) from standard tableaux with \(n\) in \(c_a\) gives
a bijection onto the standard tableaux of shape \(\mu_a\), with
inverse given by placing \(n\) in \(c_a\).

For each \(i\) we now have the equality of vector subspaces
\[
  F_i=S^\lambda\cap\bigcap_{j>i}\ker D_j .
\]
One inclusion follows from the vanishing just proved. If a vector on
the right had a nonzero coefficient at a corner with index greater
than \(i\), take the largest such index \(a\) in its standard-basis
expansion. Applying \(D_a\) gives a nonzero combination of distinct
standard polytabloids of \(S^{\mu_a}\); terms from all smaller corners
vanish, and no term from a larger corner occurs. This contradicts
membership in the kernel. The same argument includes \(i=0\).
The characterization by kernels proves invariance under \(S_{n-1}\).
On \(F_i\), the map \(D_i\) has image \(S^{\mu_i}\) and kernel
\(F_{i-1}\), by its effect on the standard bases. This proves the
quotient assertion \cite[Section 2.8]{sagan}.
\end{proof}

\begin{theorem}[Branching rule]\label{res:B05}
For \(\lambda\vdash n\) with \(n\ge 1\) and \(\nu\vdash m\)
with \(m\ge 0\), one has
\[
  \begin{aligned}
  \operatorname{Res}^{S_n}_{S_{n-1}} S^\lambda
    &\simeq \bigoplus_{\mu\text{ obtained by removing one box}} S^\mu,\\
  \operatorname{Ind}^{S_{m+1}}_{S_m} S^\nu
    &\simeq \bigoplus_{\eta\text{ obtained by adding one box}} S^\eta .
  \end{aligned}
\]
The diagrams in the two sums are obtained from \(\lambda\) and
\(\nu\), respectively. Each summand occurs once.
\end{theorem}

\begin{proof}
Maschke's theorem over \(\mathbb C\) splits the restriction filtration
of Proposition~\ref{res:B04}. Two removable boxes in distinct rows
leave distinct sequences of row lengths and therefore distinct
irreducibles by Result~\ref{res:A07}. For the induction statement
use Frobenius reciprocity and Schur's lemma, in the finite-group
forms of \cite[Chapters 8--9]{algebra-core}. For each
\(\eta\vdash m+1\), the multiplicity of \(S^\eta\) in the induced
module is
\[
  \dim_{\mathbb C}\operatorname{Hom}_{S_m}
      \bigl(S^\nu,\operatorname{Res}^{S_{m+1}}_{S_m}S^\eta\bigr).
\]
The restriction formula gives dimension \(1\) precisely when
\(\nu\) results from deleting a corner of \(\eta\). Otherwise the
dimension is zero. All irreducibles are covered by
Result~\ref{res:A07}, so complete reducibility gives the stated sum
\cite[Theorem 2.8.3]{sagan}. For \(n=1\) the restriction deletes the
only box and gives the trivial \(S_0\)-module. Induction from \(S_0\)
to \(S_1\) adds that box and gives the trivial \(S_1\)-module.
\end{proof}

\begin{corollary}[Corner recurrence]\label{res:B06}
The numbers of standard tableaux satisfy
\[
   f^{\varnothing}=1,\qquad
   f^\lambda=\sum_{c\text{ removable in }\lambda}f^{\lambda\setminus c}
   \quad(\lambda\ne\varnothing).
\]
\end{corollary}

\begin{proof}
Group standard tableaux according to the position of the greatest
label. As shown in the proof of Proposition~\ref{res:B04}, that
position must be a corner, and deleting the label and filling it
back are inverse bijections. Alternatively, the recurrence for a
nonempty diagram follows by taking dimensions in the restriction
formula and using Theorem~\ref{res:B03}
\cite[Lemma 2.8.2 and Theorem 2.8.3]{sagan}.
\end{proof}

\section{A determinant and the hook product}

\begin{theorem}[Frobenius--Young dimension determinant]\label{res:B07}
Fix \(r\ge \ell(\lambda)\) and pad \(\lambda\vdash n\) with zeros
through row \(r\). For integer \(a\) put \(q(a)=1/a!\) if
\(a\ge 0\) and \(q(a)=0\) otherwise. Then
\[
  f^\lambda=n!\det\bigl[q(\lambda_i-i+j)\bigr]_{1\le i,j\le r}.
\]
The determinant of the empty matrix is \(1\).
\end{theorem}

\begin{proof}
Let \(D(\lambda)\) be the displayed determinant without the factor
\(n!\). For \(\lambda=\varnothing\) the matrix is upper triangular
with all diagonal entries \(1\). For \(\lambda\ne\varnothing\),
form \(D_i\) by reducing \(\lambda_i\) by \(1\) in the determinant
formula, with no requirement yet that the result is a partition.
If \(i<r\) and \(\lambda_i=\lambda_{i+1}\), the new \(i\)-th row
equals row \(i+1\). If \(i=r\) and \(\lambda_r=0\), the new last row
is zero. These are exactly the invalid deletions, and they give
\(D_i=0\).

For every integer \(a\), \(q(a-1)=a q(a)\), including \(a=0\) and
negative \(a\) by the zero convention. In the determinant expansion,
the signed product belonging to a permutation \(\sigma\) is
therefore multiplied in \(D_i\) by \(\lambda_i-i+\sigma(i)\).
Summing over \(i\) multiplies that product by \(n\), since
\(\sum_i\sigma(i)=\sum_i i\). It follows that
\[
    n D(\lambda)=\sum_i D_i
      =\sum_{c\text{ removable in }\lambda}D(\lambda\setminus c).
\]
Thus \(n!D(\lambda)\) and \(f^\lambda\) have the same initial value
and satisfy the recurrence of Corollary~\ref{res:B06}. Induction on \(n\)
proves the formula for every fixed \(r\) and every diagram of at
most \(r\) rows. This is a recurrence proof of the determinant
in \cite[Theorem 3.11.1]{sagan}.
\end{proof}

For a box \(c=(i,j)\) of \(\lambda\), its \emph{hook} consists of the
box itself, the boxes to its right in its row, and the boxes below it
in its column. If \(\lambda'_j\) is the height of column \(j\), its
length is
\[
    h(i,j)=\lambda_i-j+\lambda'_j-i+1.
\]

\begin{theorem}[Hook-length formula]\label{res:B08}
For every \(\lambda\vdash n\),
\[
      \dim_{\mathbb C} S^\lambda
        = f^\lambda
        = \frac{n!}{\displaystyle\prod_{c\in\lambda}h(c)} .
\]
The empty product is \(1\). In particular the diagram with one row
and the diagram with one column each give dimension one.
\end{theorem}

\begin{proof}
Keep the number \(r\) of padded rows fixed and put
\(b_i=\lambda_i+r-i\). These are strictly decreasing nonnegative
integers. Multiply row \(i\) of the determinant in
Theorem~\ref{res:B07} by \(b_i!\). The \(j\)-th entry becomes
the falling factorial
\[
    (b_i)_{r-j}=b_i(b_i-1)\cdots(b_i-r+j+1),
\]
with \((b_i)_0=1\). If \(b_i<r-j\), the product vanishes, as does
the original entry by the zero convention. Each column consists of
values of a monic polynomial of degree \(r-j\). Starting with the
columns of lowest degree, add combinations of lower degree columns
to replace each one by the corresponding power column. The result
is the Vandermonde determinant with columns of degrees
\(r-1,r-2,\ldots,0\), whose value is
\(\prod_{i<k}(b_i-b_k)\). For example, its vanishing at coincident
variables supplies all these factors, and their degree equals
the degree of the determinant. The coefficient of
\(b_1^{r-1}b_2^{r-2}\cdots b_r^0\) is \(1\) in both expressions.
Consequently
\[
  D(\lambda)=\frac{\prod_{i<k}(b_i-b_k)}{\prod_i b_i!}.
\]

We next identify the hook lengths in row \(i\). They are distinct
positive integers at most \(b_i\): moving one box right strictly
decreases \(\lambda_i-j\) and cannot increase \(\lambda'_j\), and
the greatest possible value at \(j=1\) is \(\lambda_i+r-i\).
For each \(k>i\) the integer \(b_i-b_k\) lies in the same interval,
and these excluded integers are distinct. None is a hook length
in that row. Indeed,
\[
   h(i,j)-(b_i-b_k)=\lambda_k-j+\lambda'_j-k+1.
\]
If \(\lambda_k\ge j\), then \(\lambda'_j\ge k\), so this difference
is positive. Otherwise \(\lambda'_j<k\), so it is negative. The
\(\lambda_i\) hook lengths and the \(r-i\) excluded integers are
therefore disjoint and together number \(b_i\), exhausting
\(\{1,\ldots,b_i\}\). The argument also covers a padded zero row.
Multiplying the row identities gives
\[
   \prod_{c\in\lambda}h(c)
      = \frac{\prod_i b_i!}{\prod_{i<k}(b_i-b_k)}.
\]
The determinant formula and Theorem~\ref{res:B03} now give the
claimed dimension \cite[Theorem 3.10.2]{sagan},
\cite[Theorem 4.53]{etingof}. In a single row or column the hook
lengths are \(n,n-1,\ldots,1\), as required.
\end{proof}

\begin{example}[Small diagrams and the standard representation]\label{res:B09}
For \(n\ge 2\), the sum-zero subspace of the natural point
representation of \(S_n\) is \(S^{(n-1,1)}\). In \(S_3\) this is the
two-dimensional standard representation. The \(S_4\)-modules
\(S^{(2,2)}\) and \(S^{(3,1)}\) have dimensions \(2\) and \(3\).
Their restrictions to \(S_3\) are \(S^{(2,1)}\) and
\(S^{(3)}\oplus S^{(2,1)}\), respectively.
\end{example}

\begin{proof}
For \((n-1,1)\), the bottom-row label uniquely determines a tabloid.
Thus \(M^{(n-1,1)}\) is the point representation
\(\mathbb C^n\), with basis vectors \(v_1,\ldots,v_n\).
If the bottom label is \(b\) and the top label in its column is
\(a\), the polytabloid is \(v_b-v_a\): the only nontrivial
permutation of that column interchanges \(a\) and \(b\). Every
distinct pair can occur. These differences span the sum-zero
subspace, since any vector with coordinates \(x_1,\ldots,x_n\)
summing to zero equals \(\sum_{i<n}x_i(v_i-v_n)\).
For \(n=3\), the \((2,1)\) hook lengths are \(3,1,1\), in
agreement with its dimension \(2\).

For \((2,2)\) the hooks by rows are \(3,2\) and \(2,1\).
For \((3,1)\) they are \(4,2,1\) and \(1\). The hook formula gives
\(4!/12=2\) and \(4!/8=3\). The square diagram has only the lower
right corner, with deletion shape \((2,1)\). In \((3,1)\), deleting
the top right box gives \((2,1)\) and deleting the lower box gives
\((3)\). The decompositions follow by Theorem~\ref{res:B05}.
In the opposite direction, adjoining a box to \((2,1)\) yields
\((3,1),(2,2),(2,1,1)\). The dimensions in the induced module sum
to \(3+2+3=8\), where the last \(3\) also follows by transposition
and the sign twist of Result~\ref{res:A09}. This agrees with the
index \(4\) times the original dimension \(2\).
\end{proof}
